\documentclass[11pt,a4paper]{article}
\usepackage[T1]{fontenc}    % Indispensable pour la coupure des mots en français
\usepackage[utf8]{inputenc} 
\usepackage{amsmath}
\usepackage{amssymb}
\usepackage{amsfonts}
\usepackage{stmaryrd}
\usepackage[margin=1cm]{geometry}
\usepackage{frenchmath}

\begin{document}

% --- TITRE DE L'EXERCICE (Style identique au titre 1.7 de l'image) ---
\begin{center}
    \large \textbf{1. Etude d'une matrice à diagonale nulle}
\end{center}
\vspace{0.2cm}

% --- ENCADRÉ DE L'ÉNONCÉ (Même style visuel avec source en bas à droite) ---
\noindent\framebox[\textwidth]{%
    \begin{minipage}{0.94\textwidth}
        \vspace{0.2cm}
        Soit un entier $n \geqslant 2$. On note $A$ la matrice de $\mathcal{M}_n(\mathbb{R})$ dont les coefficients diagonaux valent 0 et les autres coefficients valent 1.
        \vspace{0.2cm}
        
        \noindent $1.$ Calculer $A^2$. En déduire que $A$ est inversible et exprimer $A^{-1}$ ;
        
        \noindent $2.$ Déterminer les valeurs propres de $A$ et les espaces propres correspondants.
        \vspace{0.1cm}
        \begin{flushright}
            (\textbf{Mines-Ponts, PC})
        \end{flushright}
    \end{minipage}%
}
\vspace{0.4cm}

\end{document}